\section{Ciclo de desarrollo y definición de flujo}
El ciclo iterativo e incremental utiliza Kanban, edición 2025.5, y SSDF 1.1, NIST SP 800-218. Kanban gobierna flujo y capacidad; SSDF integra requisitos, responsabilidades y protección en el desarrollo. Los parámetros propios de Synaptix se especifican separadamente de las prácticas de las fuentes. Las entregas contractuales siguen meses 16 y 21, con sus marchas blancas, aunque existan incrementos internos anteriores. \parencites[pp.~3--5]{sd6-kanban}[sec.~2; tabla~1]{sd6-ssdf}[art.~17, pp.~12--13]{sd6-ba}

La definición \texttt{FLUJO-6} identifica unidad de valor, principio, fin, estados, límites, políticas y expectativa de servicio. Un elemento es un resultado verificable de un paquete SD7, con requisito T-12, dominio, etapa, criterios, responsable y esfuerzo. Dividirlo crea contribuciones dentro del mismo paquete y no horas adicionales. Trabajo preparado no consume WIP hasta que se inicia; trabajo bloqueado iniciado permanece en WIP. Una cola de liberación separada conserva los productos verificados que esperan autorización.

La Tabla~\ref{tab:t10-estados} define condiciones de tránsito. El inicio es la selección efectiva de capacidad para construcción y el fin es la verificación técnica con expediente completo; la aceptación contractual y el despliegue tienen estados posteriores. Así se conserva el tiempo de espera de Producción sin confundirlo con la duración de construcción.
\begingroup\setlength{\tabcolsep}{3pt}
\begin{longtable}{L{\dimexpr.19\textwidth-2\tabcolsep\relax}L{\dimexpr.39\textwidth-2\tabcolsep\relax}L{\dimexpr.42\textwidth-2\tabcolsep\relax}}
\caption{Estados y políticas del flujo de construcción}\label{tab:t10-estados}\\\hline
\EncabezadoTabla{Estado} & \EncabezadoTabla{Entrada} & \EncabezadoTabla{Salida y responsable}\\\hline\endfirsthead
\hline\EncabezadoTabla{Estado} & \EncabezadoTabla{Entrada} & \EncabezadoTabla{Salida y responsable}\\\hline\endhead
Preparado & Requisito, alcance, recepción, contratos, dependencias y datos de prueba identificados. & Funcional y Desarrollo verifican preparación; Proyecto reserva capacidad.\\\hline
Construcción & Ejecutores disponibles y predecesoras satisfechas; registrar instante de inicio. & Código, pruebas unitarias y documentación listos para revisión; Desarrollo.\\\hline
Revisión & Solicitud de integración y reportes automáticos disponibles. & Revisor distinto del autor aprueba; observaciones devuelven a construcción.\\\hline
QA & Digest, datos y configuración identificados; capacidad de Calidad reservada. & Pruebas funcionales, integración y regresión conformes; Calidad y Funcional.\\\hline
Verificado & Definición de terminado técnica cumplida y expediente íntegro. & Registrar fin de ciclo; conservar cola de liberación con su antigüedad; Calidad.\\\hline
Liberación y recepción & Ventana, ensayos, autorización y expediente T-17/T-18 completos. & SRE despliega y observa; Contraparte acepta cuando corresponda.\\\hline
\end{longtable}\FuenteTabla{Elaboración de Synaptix a partir de Kanban y las Bases Técnicas Transversales: RT-04, pp.~10--11; RT-20, p.~35.}\endgroup
Un rechazo conserva su motivo y evidencia; no se borra el inicio para reducir el tiempo medido. Defectos posteriores originan elementos enlazados, con severidad, causa y responsable. Un cambio de recepción o alcance se somete a T-9.

\section{Capacidad, cadencias y decisiones}
Por cada persona disponible se permite un elemento activo de construcción; por cada revisora, uno en revisión; y por frente habilitado, uno esperando verificación. Un resultado con varias personas consume un elemento global y reserva capacidad a cada participante. Se prohíbe abrir nueve frentes por la sola existencia de nueve dominios. El límite se expresa como $WIP_{\max}=n_e+n_r+n_f$, donde $n_e$ y $n_r$ son plazas efectivas de ejecución y revisión y $n_f$ frentes autorizados para espera; las plazas se habilitan sólo si sus agendas de T-15 permiten usarlas. El límite global no autoriza rebasar cada estado ni asignar simultáneamente construcción y revisión a una persona.

La única clase urgente es un incidente o defecto crítico. Proyecto, Seguridad y SRE documentan por qué se interrumpe trabajo, qué se suspende y qué capacidad se reasigna. El suspendido permanece contado; si la excepción supera el límite, se identifica expresamente su duración y responsable y se restablece el límite antes de seleccionar trabajo ordinario. La urgencia no elimina pruebas, permisos ni ventanas protegidas. \parencite[Actively Managing Items in a Workflow, p.~4]{sd6-kanban}

Reposición y revisión se realizan en quincenas civiles 1--15 y 16--último día. Proyecto verifica predecesoras y capacidad; Funcional ordena elementos por obligación, dependencia, riesgo y beneficio; Desarrollo selecciona trabajo realizable. La coordinación diaria trata bloqueos y su antigüedad, con una duración objetivo de quince minutos. Revisión quincenal demuestra resultados, sin sustituir aceptación; retrospectiva registra mejora, responsable, esfuerzo y comprobación. Las duraciones son decisiones de Synaptix y se imputan a los paquetes de gestión existentes.

La expectativa inicial es terminar el 85\,\% de elementos en veinte días hábiles entre inicio de construcción y verificación. Es una hipótesis operativa de Synaptix para orientar revisión de elementos envejecidos, no un percentil medido ni garantía contractual. Se recalibra con el percentil 85 de tiempos completos observados por clase; el reporte conserva muestra, dispersión y bloqueados abiertos. La previsión no reemplaza el plazo de cada paquete ni justifica una selección que ya exceda capacidad. \parencite[Definition of Workflow y Flow Metrics, pp.~3--5]{sd6-kanban}

\section{Requisitos, diseño, código y dependencias}
La ficha \texttt{ITEM-6} conserva ID, T-12 y complemento del caso, paquete, dominio, actor, contrato, criterios positivos/negativos, etapa, esfuerzo, riesgo, responsable, revisor, versión y estados. Cambiar una implementación no permite perder una cláusula del requisito. Las pruebas comprueban autorización, duplicados, errores de dependencia, desconexión, conflicto y conciliación, además del recorrido exitoso.

El registro de decisión de arquitectura identifica problema, alternativas, elección, razones, impacto y prueba. Una modificación relevante actualiza modelo de amenazas por componente/integración, contratos e interfaces; el Comité de Arquitectura la revisa. Se conservan autoridades cloud/local y continuidad de SD3--SD5. La evolución favorece cambios compatibles y refactorización verificable, con recepción objetiva y presupuesto de horas. \parencites[RT-11.02, p.~23]{sd6-bt}[PW.1/PW.2, tabla~1]{sd6-ssdf}

La política de protección de ramas, versión 1.0, se aplicará así. Las ramas \texttt{main} (principal) y \texttt{release/*} estarán protegidas: se deniegan escritura directa, force-push, eliminación y bypass, incluso a administradores y cuentas de servicio. Sólo se integra mediante solicitud de revisión, aprobación de una persona distinta del autor y controles obligatorios conformes. La aprobación queda ligada al commit revisado; un commit nuevo invalida la aprobación y repite controles. La urgencia conserva estas puertas. El responsable del repositorio verificará permisos efectivos con intentos negativos antes de habilitarlo. Las ramas cortas se enlazan al elemento, y las versiones se identifican por commit y digest inmutables. El estándar de Java/Kotlin y TypeScript/JavaScript conserva contratos tipados, validación de entradas, permisos en servidor, transacciones e idempotencia, errores explícitos y auditoría sin secretos. Una excepción al estándar documenta motivo y prueba; no permite debilitar las bases. \parencites[RT-04.03/04, p.~10]{sd6-bt}[PS.1; PW.5/PW.7, tabla~1]{sd6-ssdf}

Antes de adoptar dependencias, Desarrollo registra componente, origen oficial, versión fija, licencia, soporte, vulnerabilidades y alternativa. Seguridad y Arquitectura aprueban; un componente con vulnerabilidad conocida no remediada no entra al conjunto aprobado. SBOM identifica las dependencias efectivas de la versión. Las actualizaciones vuelven a pasar análisis y regresión; los artefactos de herramientas del pipeline también tienen origen, huella y revisión. \parencites[RT-11.23/26, p.~24]{sd6-bt}[PW.4, tabla~1]{sd6-ssdf}

El manifiesto \texttt{LIB-6}, versión 1.0, enlaza cada requisito T-12 y complemento aplicable con la incidencia \texttt{ITEM-6}, solicitud de revisión, aprobación, commit integrado, contratos y esquema/migración versionados, casos y ejecuciones de prueba, reportes, SBOM, firma y digest del artefacto. Cada ejecución identifica resultado observado, instante, datos y ambiente; una prueba sólo planificada no satisface el enlace. El registro de despliegue añade ID de ejecución del pipeline, configuración/IaC, ambiente destino, digest desplegado, autorizador, inicio/fin, resultado y comprobación de salud; si hay retorno, conserva ejecución y digest recuperado. Antes de desplegar se bloquea cualquier enlace previo ausente, divergente o fallido. Después se incorpora la evidencia efectiva de despliegue antes de cerrar la incidencia; un fallo conserva el estado fallido y el retorno, sin presentarse como entrega exitosa. Los registros son inmutables y las correcciones añaden una versión enlazada. Calidad comprobará el recorrido en ambos sentidos, desde requisito a despliegue y desde digest a requisitos, incluyendo una promoción rechazada por eslabón faltante.

\section{Ambientes y controles DevSecOps}
Desarrollo, QA, Preproducción, Producción y DR estarán operativos en H3 de E-25. Cada uno tiene acceso, configuración, secretos y datos propios. Preproducción reproduce topología, versiones y controles; las diferencias de capacidad se declaran sin eliminar redundancia. Se usan datos sintéticos o protección verificable, con reinicio de pruebas a estado conocido. OpenTofu conserva versiones de infraestructura y estado cifrado. \parencites[RT-04.01/02, p.~10; RT-11.25, p.~24]{sd6-bt}[PO.5, tabla~1]{sd6-ssdf}

Para recibir H3, SRE y Calidad comprobarán en cada ambiente acceso nominativo, aislamiento, inventario/configuración IaC, secretos propios y ejecución de una versión identificada. Desarrollo se reconstruirá desde código con datos sintéticos o anonimizados; QA conservará datos controlados/versionados y demostrará reinicio al estado conocido. Preproducción tendrá la misma topología, configuración y versiones que Producción, con datos representativos protegidos; sólo variará la capacidad declarada en SD4. Producción mantendrá acceso restringido y auditado, sin acceso interactivo directo de desarrolladores. DR conservará replicación continua y ejecutará conmutación y retorno reales al habilitarse y semestralmente, dentro de RTO/RPO de SD4. El acta H3 identificará ambiente, versión, comprobaciones, resultados y responsables; un ambiente ausente o una comprobación fallida impide recibir el hito. Son condiciones de ejecución futura, no resultados ya obtenidos.

La Tabla~\ref{tab:t10-pipeline} especifica el flujo de SD4. AWS CodeBuild y AWS CodePipeline V2 orquestan; T-11, ítems 20A--20F, individualiza Semgrep 1.180.0, Trivy 0.75.0, Gitleaks 8.30.1, ZAP 2.17.0 y Cosign 3.1.3. OpenTofu 1.13.1 corresponde al ítem 20. Cada ejecución conserva la configuración de los servicios administrados, versión/digest de imagen de construcción, herramientas, reglas y bases de vulnerabilidades; Seguridad comprobará vigencia y aprobará cualquier actualización antes de usarla. El mismo artefacto se promueve por digest sin recompilación, y configuración/secretos se inyectan por ambiente. \parencite[RT-04.05--11, pp.~10--11; RT-11.22--24, p.~24]{sd6-bt}
\begingroup\setlength{\tabcolsep}{3pt}
\begin{longtable}{L{\dimexpr.20\textwidth-2\tabcolsep\relax}L{\dimexpr.48\textwidth-2\tabcolsep\relax}L{\dimexpr.32\textwidth-2\tabcolsep\relax}}
\caption{Controles bloqueantes y evidencias del pipeline}\label{tab:t10-pipeline}\\\hline
\EncabezadoTabla{Paso} & \EncabezadoTabla{Bloqueo} & \EncabezadoTabla{Producto y responsable}\\\hline\endfirsthead
\hline\EncabezadoTabla{Paso} & \EncabezadoTabla{Bloqueo} & \EncabezadoTabla{Producto y responsable}\\\hline\endhead
Trazabilidad & Sin requisito, elemento, commit resuelto o revisión por pares. & Manifiesto; Desarrollo.\\\hline
Construcción & Compilación o prueba unitaria fallida, reporte ausente o cobertura de negocio menor a 70\,\%. & Reportes y digest; Desarrollo/Calidad.\\\hline
SAST/SCA & Hallazgo crítico/alto abierto, dependencia sin aprobación o licencia inadmisible. & Reportes Semgrep/Trivy y decisión; Seguridad.\\\hline
Secretos e imágenes & Secreto detectado, imagen sin análisis, hallazgo crítico/alto abierto en imagen o infraestructura con control incumplido. & Gitleaks/Trivy; Seguridad/SRE.\\\hline
DAST & Falla autenticada, autorización defectuosa o hallazgo crítico/alto. & ZAP y pruebas adversas; Seguridad/Calidad.\\\hline
Firma y procedencia & SBOM ausente o firma, builder, fuente, parámetros o digest no verificables. & Atestación y SBOM; SRE/Seguridad.\\\hline
Promoción & Pruebas, intrusión independiente, retorno, ventana o autorización incompletos. & Expediente de liberación; Proyecto/SRE.\\\hline
\end{longtable}\FuenteTabla{Bases Técnicas Transversales: RT-04, pp.~10--11; RT-11, pp.~23--24; cap.~20, pp.~34--35. Aplicación al pipeline de SD4.}\endgroup
La cobertura divide líneas ejecutables de negocio ejercitadas por el total, sin excluir funciones para alterar el denominador. Toda prueba conserva versión, ambiente, datos, instante, resultado y evidencia. La herramienta y la base de análisis quedan identificadas; un reporte de otra versión no habilita el digest actual. \parencite[PO.3/PO.4; PW.7/PW.8, tabla~1]{sd6-ssdf}

SLSA Build L3 versión 1.2 exige comprobar aislamiento de ejecuciones y procedencia generada por un plano de control separado, que el código del build no pueda falsificar. La identidad de firma no se expone a los pasos de construcción; el verificador comprueba digest, fuente, parámetros y plataforma autorizada. No se atribuye automáticamente L3 al servicio administrado ni a la existencia de una firma. El expediente de habilitación contiene pruebas negativas de esos controles. \parencites{sd6-slsa,sd6-slsa-requisitos}

Desarrolladores no tienen acceso interactivo directo a Producción. La excepción requiere incidencia, aprobación, MFA, sesión grabada y vencimiento; se conserva el máximo de sesenta minutos del SD4. Si la modalidad no permite grabar, se deniega y se usa recuperación autorizada. La separación de funciones permanece durante contingencia. \parencite[RT-11.27, p.~24]{sd6-bt}

La política SEC-ROT-1.0 de SD4 concreta Secrets Manager/Lambda y GLS-1.0 local: rotación automática cada treinta días para credenciales persistentes, prueba y revocación, transición máxima de quince minutos, auditoría de cada acceso/resultado y prohibición de credenciales embebidas. La renovación de credenciales temporales respeta su expiración. La rotación local continúa sin WAN; una credencial externa vencida no habilita su integración. Calidad y Seguridad ensayan el ciclo y las fallas antes de habilitar el gestor o la versión; LIB-6 registra evidencia sin secretos. \parencite[RT-04.09, p.~11]{sd6-bt}

\section{Pruebas, definición de terminado y aceptación}
La definición de terminado se acordará con el CLIENTE para cada entregable. Su verificación exige requisitos cubiertos, revisión de código, pruebas conformes, datos conciliados, seguridad, documentación y despliegue aplicable. Un incremento interno termina técnicamente tras QA y expediente; la recepción contractual añade ensayos, autorización y acta. El estado separado evita declarar aceptado lo que sólo está construido. \parencite[RT-20.07/08, p.~35]{sd6-bt}

Unitarias/componentes corren continuamente, con cobertura de negocio de al menos 70\,\% y ninguna falla; integración tras cada despliegue QA ejecuta todos los flujos aplicables sin error; regresión automatizada completa precede a cada promoción a Preproducción; UAT precede a cada certificación de etapa y exige firma de la Contraparte. Antes de cada paso a producción se ejecutan carga y estrés a 1,5 veces el peak declarado, accesibilidad WCAG 2.2 AA automática y manual documentada, seguridad ofensiva independiente sin críticos/altos abiertos, resiliencia con recuperación sin intervención y conmutación real DR dentro de RTO/RPO. Seguridad ofensiva se repite además anualmente, y resiliencia/DR al menos semestralmente. El paquete de liberación identifica umbral, volumen, mezcla, digest y evidencia efectiva; una prueba de otra versión no habilita el artefacto nuevo. Calidad consolida reportes sin duplicar esfuerzo de construcción. Funcional y Contraparte validan procesos; Seguridad recibe informe íntegro de intrusión y remediación; SRE mide retorno, RTO/RPO y continuidad. \parencite[RT-11.20, p.~24; cap.~20, pp.~34--35]{sd6-bt}

El procedimiento PRUEBAS-6 aplica ISO/IEC/IEEE 29119-2:2021 para procesos, 29119-3:2021 para documentación y 29119-4:2021 para técnicas. Calidad mantiene la correspondencia entre proceso, control, registro y criterio; las fichas de SD7 reservan su ejecución y SD9/T-13/T-17 organizan calidad y recepción. La Tabla~\ref{tab:t10-29119} concreta la aplicación, sin atribuir certificación ni resultados ejecutados. \parencites[art.~4.3, p.~5]{sd6-ba}{sd6-29119-2,sd6-29119-3,sd6-29119-4}
\begin{table}[htbp]
\caption{Aplicación del marco de pruebas a una versión}\label{tab:t10-29119}
\begin{tabularx}{\linewidth}{L{.21\linewidth}Y L{.25\linewidth}}
\hline
\EncabezadoTabla{Proceso} & \EncabezadoTabla{Control aplicado} & \EncabezadoTabla{Registro y responsable}\\\hline
Gobierno y planificación & Alcance T-12/CA, tipos, riesgos, ambientes, datos, capacidad y entradas/salidas; aprobar antes del ensayo. & Plan PRUEBAS-6; Calidad y Proyecto.\\\hline
Diseño e implementación & Particiones y límites para cantidades/tiempos; tablas de decisión para permisos; estados para salida, regreso, retiro y retorno. & Caso con precondición, estímulo y resultado esperado; Calidad/Funcional.\\\hline
Entorno y ejecución & Configuración, digest y datos fijos; registrar secuencia, resultado observado y evidencia; impedir promoción ante falla. & Bitácora por ejecución; Calidad y SRE.\\\hline
Seguimiento e incidencias & Comparar cobertura y criterios; clasificar defectos, corregir y repetir prueba y regresión afectadas. & Reporte de avance/defecto enlazado; Calidad y Desarrollo.\\\hline
Cierre & Confirmar criterios, diferencias resueltas y evidencia; declarar pendientes y conservar informe para recepción. & Informe de cierre y expediente T-17; Calidad y Contraparte.\\\hline
\end{tabularx}
\FuenteTabla{Aplicación de Synaptix a los procesos, documentación y técnicas de ISO/IEC/IEEE 29119-2/3/4:2021.}
\end{table}
Cada caso conserva ID de requisito, versión, técnica, precondición, datos, pasos, resultado esperado/observado, ejecutor y referencia de defecto. Calidad verifica cobertura de requisitos y riesgos antes de cerrar; usar una técnica no reduce el catálogo obligatorio del capítulo 20. El reporte diferencia plan, ejecución, aprobación técnica y aceptación contractual.

Las pruebas específicas cubren 24 horas de autonomía del recinto, 12 de dispositivos de terreno y 8 de embarcación de inspección, con registros íntegros y sincronización conforme a SD4/SD5. Incluyen salidas/regresos, escuela y retiro de menores, expendio de combustible, amarras y notificación por marejada. El ensayo conserva variantes y rechazo de operaciones sin permiso, no sólo conectividad del equipo. \parencite[cap.~15, RT-03.10, p.~33]{sd6-c07}

Los dos ensayos de migración por etapa y la marcha blanca siguen SD7/T-18, con conciliación diaria y ninguna diferencia no explicada. Se evalúan indicadores y umbrales de cierre de cuatro semanas según el artículo 17.3. El acta precede al cambio oficial de autoridad; la sesión de revisión y contingencia siguen el supuesto ya desarrollado en SD7. Una observación impeditiva bloquea el paso y no se sustituye por aceptación tácita. \parencites[art.~17.3, pp.~12--13]{sd6-ba}[RT-20.02--04, p.~35]{sd6-bt}

\section{Liberación y retorno}
Cada liberación identifica servicio, requisito, commit, esquema, digest, SBOM, firma/procedencia, reportes, datos, responsables, ventana, plan y criterio de retorno. SRE deposita antes de promover las fuentes, configuración, IaC, scripts y manifiesto reproducible del commit en el repositorio del CLIENTE; Calidad comprueba versión, integridad, acceso y restauración del paquete. El pipeline bloquea si falta ese depósito o la evidencia obligatoria. El responsable de custodia de SD7 prepara además actualización semestral ante tercero independiente, con comprobante y condiciones de liberación por insolvencia o incumplimiento grave; ese depósito adicional no reemplaza la actualización de cada liberación. \parencite[art.~84.1/6, p.~43]{sd6-ba}

\textbf{Orquestación de liberación DEP-6 1.0.}
DEP-6 es el procedimiento versionado de despliegue y retorno, ejecutado por CodePipeline/CodeBuild en nube y por el agente autorizado del Runtime en el recinto. El manifiesto LIB-6 fija digest nuevo y anterior, configuración por destino, plan MIG-6, comprobaciones, umbrales, orden, contratos y plan de retorno. La autorización de cambio precede al despacho del job; a partir de ella, preparación, ejecución, comprobación y reversión son automáticas, sin comandos humanos, consola, SSH ni copia manual de archivos en Producción. El agente sólo acepta planes firmados de la identidad autorizada, con destino, vigencia y secuencia válidos; obtiene secretos en memoria desde GLS-1.0, sin credenciales en argumentos o registros. El paquete completo y el anterior se conservan en ambos nodos antes de comenzar, permitiendo ejecutar el retorno local aunque se pierda la WAN. Una nueva liberación no empieza sin autorización y expediente completos.

La máquina de estados persistente registra preparado, validado, migración compatible, candidato sano, promoción, observación, consolidado y retorno. Un ID único identifica cada ejecución; cada paso comprueba precondición y poscondición y registra huella, destino, instante y resultado. Tras reiniciar, el agente consulta estado efectivo, autoridad HA y diario antes de reanudar el mismo paso idempotente. Un bloqueo exclusivo impide dos liberaciones o migraciones simultáneas del mismo destino; los trabajadores usan token de autoridad y cesan si pierden el rol. DEP-6 valida Pacemaker/Corosync, fencing y réplica síncrona antes de alterar el par: no crea un segundo escritor PostgreSQL ni promueve un nodo cuyo estado sea incierto. Firma, paquete, réplica, capacidad o condición de salud divergentes detienen el avance. Una falla del candidato devuelve automáticamente el tráfico al digest anterior, cancela nuevas transformaciones y aplica, cuando corresponda, la reversa MIG-6 ya ensayada; verifica servicio e integridad antes de cerrar el retorno. Si no existe autoridad segura, conserva evidencias, bloquea escrituras inseguras y escala incidente, sin declarar recuperación conseguida. \parencite[RT-04.06, p.~11]{sd6-bt}

\textbf{Azul-verde de aplicación y drenaje.}
En nube se mantienen conjuntos N y N+1 sanos, en ambas zonas, detrás del balanceador existente. Se enruta 10\,\% del tráfico al candidato durante cinco minutos; sólo con comprobaciones conformes pasa a 100\,\% y observa quince minutos. Una comprobación crítica fallida, SLO del flujo incumplido o integridad fallida activa retorno automático. Los porcentajes y períodos son decisiones de diseño, no resultados medidos. El grupo anterior permanece disponible durante observación y ventana de retorno. No se detiene una tarea con peticiones en curso: primero se retira de admisión, se drena y se comprueba el contador de trabajo pendiente. Un drenaje vencido bloquea retirada y conserva el grupo anterior, en lugar de matar conexiones para terminar la liberación.

Para cambios compatibles de la aplicación local se seleccionan dos slots Podman N/N+1 sobre el nodo activo, detrás de un módulo de encaminamiento del Runtime con ciclo de vida separado del ejecutable de negocio. Este módulo será desarrollo propio de Synaptix: conserva el ingreso HTTPS y el mapa de destinos, cambia ese mapa atómicamente y mantiene conexiones en curso con el slot anterior. No se instala un nuevo producto de tercero ni se atribuye esta función a Pacemaker. Primero se prepara y verifica el slot nuevo en el pasivo sin tráfico de escritura; después se prepara en el activo, se verifica contra la autoridad vigente y se aplica la misma pauta 10\,\%/cinco minutos/100\,\%/quince minutos. PostgreSQL 18, su primario, IP virtual, GLS y el broker conservan su rol durante esa promoción: desplegar la aplicación no exige conmutar la base ni reiniciar todo el nodo. El pasivo conserva ambos digests y el mapa consolidado para contingencia, sin ejecutar dos escritores. Ambos slots y el módulo se incluyen en el ensayo de capacidad EQ-4; CPU, memoria, almacenamiento e IOPS del nodo permanecen como máximo al 60\,\% de la capacidad definida. Si el solapamiento incumple este límite, se redimensiona y ensaya antes de promover, sin consumir el 40\,\% de reserva ni detener N para alojar N+1.

Las sesiones y la idempotencia se conservan fuera del slot en PostgreSQL; N/N+1 usan contratos y esquema compatibles conforme a MIG-6. Una operación de escritura lleva ID idempotente y resultado durable confirmado transaccionalmente; un reintento recupera ese resultado y no repite despacho, cargo ni notificación externa. Los workers dejan de tomar trabajo antes del drenaje y concluyen o devuelven al diario/cola el trabajo no confirmado. Una conexión larga mantiene su slot hasta cierre, o migra mediante aviso, cursor confirmado y reconexión automática ensayada; nunca se interpreta un timeout como confirmación de negocio. Los reintentos usan demora acotada, retroceso y jitter dentro del plazo del flujo. Una conexión no migrable conserva N y bloquea retirada; un reintento que exceda el plazo contractual cuenta como incumplimiento aunque al final recupere el registro.

Antes de cada paso a Producción, Calidad y SRE ejecutarán en Preproducción EQ-4 el plan exacto DEP-6 con el mismo digest, contratos y configuración normalizada, a 1,5 veces el peak habilitado. El ensayo cubre promoción y retorno con tráfico continuo cloud/local, peticiones en vuelo, conexiones largas, repetición de escrituras, pérdida WAN, reinicio del agente y falla inducida del candidato. LIB-6 conserva línea de tiempo de estados/tráfico, contadores, latencia por flujo, errores visibles, IDs y conciliación antes/después; sólo habilita producción con cero solicitudes de negocio perdidas o duplicadas, cero interrupciones del servicio y plazos de cada flujo satisfechos. La ausencia de error en la sonda de salud no reemplaza esta comprobación de extremo a extremo. Se libera por proceso o zona del Caso 07, con sus dependencias y ventanas protegidas. \parencites[RT-04.07, p.~11]{sd6-bt}[art.~20, pp.~14--15]{sd6-ba}[cap.~13, numeral 13.3, pp.~29--30]{sd6-c07}

\textbf{Límite del camino de continuidad.}
El camino azul-verde anterior cubre ejecutables y configuración compatibles de aplicación; no acredita liberación sin interrupción para sustituir el módulo de encaminamiento, reiniciar el sistema operativo o conmutar el primario PostgreSQL, la IP virtual, GLS o EMQX. Para esas clases, DEP-6 automatiza actualización secuencial, drenaje, conmutación bajo HA/fencing, comprobación y retorno, con réplica sana y paquetes previos; nunca actualiza ambos nodos a la vez. La conmutación puede cortar conexiones y el objetivo de recuperación local de SD4 no equivale a cero interrupción. Esas clases no se autorizan como liberación conforme a RT-04.07 hasta disponer de un camino compatible y ensayado de continuidad de extremo a extremo. Programar una ventana o conservar registros mediante reintentos no satisface por sí solo la cláusula. El expediente conserva esta brecha de diseño y su resolución como condición de promoción, sin presentar la prueba futura como resultado obtenido.

\textbf{Migración automatizada de esquema MIG-6 1.0.}
Cada liberación incluye scripts SQL de avance y reversa, pruebas y pre/postcondiciones, orden, versión anterior/destino y huellas inmutables, vinculados al commit y digest de LIB-6. El ejecutor MIG-6 1.0 utiliza PostgreSQL 18 y su cliente de la imagen aprobada; CodeBuild ejecuta el plan cloud y el agente autorizado del Runtime ejecuta el mismo plan firmado en el par local. La aprobación habilita el job: no sustituye su ejecución por comandos manuales. Sólo actúa el escritor autorizado, con exclusión mediante bloqueo asesor y comprobación de HA/fencing; la identidad de migración es distinta de la aplicación y obtiene su credencial en memoria desde el gestor. La tabla de historial registra ID, huella, versión, destino, inicio/fin y resultado. Una huella modificada, versión de partida divergente o paso fallido impide continuar; la corrección crea una migración nueva y preserva el historial.

El avance usa incorporación, transición y retirada. Primero crea estructuras compatibles; después transforma en lotes idempotentes con correspondencia de claves y conciliación; sólo retira tras vencer la ventana de retorno aprobada y comprobar la reversa. Durante el despliegue, las versiones N y N+1 leen/escriben mediante contratos compatibles, vistas/adaptadores y reglas de duplicación controladas por la aplicación. Se mantiene la representación anterior y un diario ordenado de cambios, con identificador idempotente, antes/después y esquema de origen; los cambios se confirman transaccionalmente con ese diario. Una función exclusiva N+1 permanece desactivada hasta superar la transición; no se admiten escrituras que la reversa no pueda representar.

Antes de cada cambio se verifica respaldo recuperable y la conservación de originales/diario de retorno. El pipeline ensaya avance, ambas versiones concurrentes, interrupción/reanudación y reversa automática con hechos nuevos aceptados durante la ventana. Ante precondición, conciliación o integridad fallida, MIG-6 detiene nuevas transformaciones, ejecuta los pasos inversos ensayados y concilia los hechos posteriores sin repetir efectos ni perder operaciones. El cambio no se declara recibido hasta comprobar autoridad, restricciones, conteos y totales de control. Un paso no transaccional declara su compensación y checkpoints explícitos. Una retirada destructiva sin reversa ensayada y datos suficientes se bloquea; una copia anterior aislada no basta porque omitiría los hechos nuevos. Si la reversa directa exige restauración, el plan automatiza restauración aislada, reaplicación del diario y validación antes de conmutar la autoridad, conservando el origen y los límites de continuidad aplicables. No sobrescribe a ciegas la base activa ni repite llamadas externas confirmadas.

Desarrollo entrega scripts y contratos; SRE ejecuta mediante pipeline; Calidad verifica compatibilidad, avance/retorno y conciliación cloud/local. LIB-6 conserva reportes, huellas y registro de ejecución real; no habilita una migración basada sólo en una prueba prevista. Este control corresponde al esquema de cada versión y no sustituye los dos ensayos de migración del legado de SD5/SD7. \parencite[RT-04.10, p.~11]{sd6-bt}

El esquema PostgreSQL usa incorporación, transición y retirada. Se ensayan dos versiones coexistentes, lotes idempotentes y conciliación; los originales permanecen hasta cerrar retorno. Un retiro destructivo requiere copia verificable, restauración y reaplicación ensayadas. El rollback de código no se confunde con recuperación de datos. \parencite[RT-04.10, p.~11]{sd6-bt}

SRE comprueba ventana, autorización y aviso recibido por el CLIENTE con al menos diez días hábiles antes de un mantenimiento programado. El expediente conserva envío, destinatario, alcance, impacto, fecha y retorno; una reprogramación comprueba de nuevo el preaviso. El congelamiento total, campañas de varadero y eventos náuticos bloquean una liberación no permitida. Ante aviso de marejada se suspende inmediatamente toda actividad programada del proyecto, incluido trabajo aislado o remoto; la reanudación requiere cese de la condición y nueva comprobación de ventana y capacidad. La restauración de un servicio existente se rige por el procedimiento de incidente y su autoridad, sin usar la urgencia para introducir una evolución no aprobada. Las actividades necesarias de recuperación preservan seguridad de personas y continuidad. \parencites[cap.~10, restricción 12, p.~26; numeral 13.2, p.~29; cap.~15, RT-10.05, p.~34]{sd6-c07}[art.~20, p.~15]{sd6-ba}[RT-10.05, p.~22]{sd6-bt}

\section{Métricas, deuda y mejora}
El reporte diario de flujo recoge WIP, terminados por período, antigüedad de abiertos y tiempo de ciclo, con unidades y extremos definidos en \texttt{FLUJO-6}. Se conservan bloqueos y reaperturas. La revisión quincenal analiza esperas, defectos y variación por clase, y selecciona mejoras dentro de capacidad. \parencite[Flow Metrics, p.~5]{sd6-kanban}

Las métricas de RT-04.12 tienen estas definiciones y objetivos propios de Synaptix:
\begin{itemize}
\item Frecuencia: liberaciones exitosas por servicio y mes. Objetivo de al menos una en cada mes con ventana y cambio autorizado; sin trabajo autorizado, registrar cero y su causa.
\item Tiempo de entrega: desde commit integrado hasta puesta efectiva en Producción. Objetivo de hasta veinte días hábiles para cambios ordinarios que disponen de ventana; el indicador observado incluye toda espera y se reporta por separado su causa de calendario.
\item Tasa de cambios fallidos: despliegues que requieren reversión, corrección urgente o causan incidente, divididos por despliegues productivos del período. Objetivo inferior al 5\,\%; sin despliegues no se declara tasa cero, sino no calculable.
\item Restauración: tiempo desde detección del fallo atribuible al despliegue hasta recuperar el servicio y comprobar integridad. Objetivo de hasta sesenta minutos para reversión de software, respetando cualquier plazo contractual más estricto; recuperación de datos y DR conservan los RTO/RPO de SD4 y se reportan separadamente.
\end{itemize}
Estos objetivos no son valores de catálogo, resultados observados ni excepciones al calendario contractual. SRE obtiene instantes del pipeline y monitorización; Calidad comprueba población y cálculo. La frecuencia de cambios nunca prima sobre una puerta de seguridad o el presupuesto de error del servicio. \parencite[RT-04.12, p.~11; RT-10.09, p.~22]{sd6-bt}

La deuda registra componente, causa, riesgo, HH de eliminación, prioridad, responsable, pruebas y cierre. Se mantiene la reserva F24 de 32 HH mensuales de SD7: $32/240=0,1333$, equivalente a 13,33\,\% de las cuatro bolsas de ingeniería de Operación; la reserva interna de mantenimiento es $24\times0,20=4,8$ HH. La capacidad de turnos no entra en ese denominador. Una acción cargada a F24 no se carga de nuevo al mantenimiento. Desarrollo informa horas consumidas y riesgo/HH realmente reducidos; cerrar un ticket sin prueba no demuestra reducción. \parencite[RT-21.21, p.~38]{sd6-bt}

Seguridad realizará evaluación inicial y anual OWASP SAMM, conservando evidencias y acciones por práctica conforme al diseño SD4. SSDF orienta además identificar vulnerabilidades, priorizar su resolución y analizar causas para prevenir recurrencia. SRE aporta incidentes; Desarrollo incorpora correcciones; Calidad comprueba regresión y cierre. El CLIENTE recibe expediente de versión y documentación mantenible, sin atribuir una certificación SAMM o SSDF obtenida. \parencites{sd6-samm}[RV.1--RV.3, tabla~1]{sd6-ssdf}
